% ---------
% --------
%!TEX TS-program = pdflatex
%!TEX encoding = UTF-8 Unicode
%!TEX spellcheck = English (Aspell)

\documentclass[11pt]{article}
\usepackage{lucy,handout,graphicx}
\usepackage{fourier-orns}
\usepackage{mathtools}
\usepackage[normalem]{ulem} % either use this (simple) or
\usepackage{tikz}

\def\Cdot{\mathbin{\text{\normalfont \textbullet}}}
\def\Sym#1{\texttt{\color{BrickRed}#1}}
\def\BlueSym#1{\textbf{\texttt{\color{RoyalBlue}#1}}}
\allowdisplaybreaks

\begin{document}

\begin{center}
\Large\textbf{CSCI 332, Fall 2026}%
\\
\LARGE\textbf{Homework 9 (Part 2---Part 1 is on Prairie Learn)}%
\\[0.5ex]
\large Due Monday, November 2 Anywhere on Earth (6am Tuesday)
\end{center}

\bigskip
\hrule
\bigskip

\subsection*{Submission Requirements}
\begin{itemize}
    \item Type or clearly hand-write your solutions into a PDF format so that they are legible and professional. Submit your PDF on Gradescope. 
    \item Do not submit your first draft. Type or clearly re-write your solutions for your final submission. If your submission is not legible, we will ask you to resubmit.
    \item Use Gradescope to assign problems to the correct page(s) in your solution. If you do not do this correctly, we will ask you to resubmit.
\end{itemize}

\subsection*{Academic Integrity}

Remember, you may access \EMPH{any} resource in preparing your solution to the homework. However, you \EMPH{must}
\begin{itemize}
    \item write your solutions in your own words, and
    \item credit every resource you use (for example: ``Bob Smith helped me on
    this problem. He took this course at UM in Fall 2020''; ``I found a solution
    to a problem similar to this one in the lecture notes for a different
    course, found at this link: www.profzeno.com/agreatclass/lecture10''; ``Here is a link to the full transcript of my chat with ChatGPT about the problems on this HW''.)
    Ask if you need any clarifications about how to properly cite your sources!
\end{itemize}




\newpage
%----------------------------------------------------------------------
    
\headers{CSCI 332}{Homework 9 (due November 2)}{Fall 2026}


\begin{enumerate}
%\parindent 1.5em \itemsep 4ex plus 0.5fil

%----------------------------------------------------------------------

\item (0 points) List the resources you used for each problem in this homework. If you used no resources except
lecture notes and videos and the linked textbook, you can say ``none''. If you used
an AI tool, you should put a link to your full conversation here, and you must
use the
\href{https://umaal.org/files/teaching/csci332_fall2026/agents.md}{provided
prompt}
 to guide the AI to act
as a tutor, not an answer-provider.

You will be asked to resubmit your homework if this section is blank.

Remember that no matter what resources you consulted, you must write your solutions
\emph{in your own words}.

\item Your best friend is attending college on the other side of the country. They have
a BIG secret to share, but they don't trust any communication method except face-to-face. (What if the phones
are tapped?) Thus, (blowing off all of your responsibilities), you have decided that you will
both get in your cars and start driving toward one another. However,
(to minimize the impact of said blowing off of responsibilities), you want to meet in the city in the middle
that you can both reach as soon as possible.

You have the map of the country stored as a weighted graph $G=(V, E)$
where the vertices $V$ are the cities and the edges $E$ are roads
that directly connect two cities. Each edge has a weight $t(e)$
equal to the time it drive between the edge's endpoints.
One vertex is labeled $y$, where you live, and another is labeled $f$,
where your friend lives.

Your goal will be to develop an algorithm to find the city $c$
where you and your friend should meet.

\begin{enumerate}
  \item For each of the following algorithms, give a counterexample
  (an example $G$ and the corresponding $c$) showing that the algorithm
  \emph{does not work}.

  \begin{enumerate}
    \item Find the mid point of the shortest path from $y$ to $f$.
    \item Find the mid point of the path from $y$ to $f$ in the minimum spanning tree of $G$.
    \item Find the vertex $c$ that minimizes the sum of the shortest path
    distance from $y$ to $c$ and the shortest path distance from $f$ to $c$.
  \end{enumerate}

  \item Give a (correct) algorithm to find $c$ and analyze its runtime.

\end{enumerate}






\end{enumerate}
%%%%%%%%%%%%%%%%%%%%
\end{document}
